% ECONOMICS_PROBLEM: {"id": "E62-411", "title": "Targeted transfers and automatic stabilizers", "jel_code": "E62", "source_id": "OP-0253"}
% !TeX program = xelatex
\documentclass[11pt,a4paper]{article}
\usepackage[margin=23mm]{geometry}
\usepackage{fontspec}
\setmainfont{Latin Modern Roman}
\usepackage{amsmath,amssymb,mathtools}
\usepackage{xcolor,enumitem,booktabs,longtable,array,xurl,hyperref}
\definecolor{muted}{HTML}{53636C}
\hypersetup{colorlinks=true,urlcolor=blue,linkcolor=blue}
\setlength{\parindent}{0pt}
\setlength{\parskip}{5pt}
\newcommand{\R}{\mathbb R}
\newcommand{\E}{\mathbb E}
\newcommand{\N}{\mathbb N}
\newcommand{\ind}{\mathbf 1}
\newcommand{\Var}{\operatorname{Var}}
\newcommand{\Cov}{\operatorname{Cov}}
\newcommand{\supp}{\operatorname{supp}}
\DeclareMathOperator*{\argmin}{arg\,min}
\DeclareMathOperator*{\argmax}{arg\,max}
\newcommand{\entrymeta}[3]{{\sffamily\small\color{muted}\textbf{\texttt{#1}}\quad|\quad #2\par #3\par}}
\newcommand{\Problemref}[1]{\texttt{#1}}
\newcommand{\JELref}[2]{\texttt{#2} (JEL \texttt{#1})}
\begin{document}
\section*{Targeted transfers and automatic stabilizers}
\textbf{Problem ID:} \texttt{E62-411}\quad \textbf{JEL:} \texttt{E62}\par
% BEGIN ECONOMICS STATEMENT
\label{problem:OP-0253}
{\sffamily\small\color{muted}Original subject group: Fiscal policy and sovereign debt\par}
\label{definitions:problem:30007534}
\textbf{Audit classification:} Published research agenda. Optimal heterogeneous-agent policy is an agenda; exact transfer targeting under administrative constraints is editorial. Verification concerns the cited version, not first priority or proof that this exact editorial specification remains unsolved on October 8, 2026.
\subsection*{Definitions and precise research target}
\paragraph{Editorial mathematical specification.} Fix a heterogeneous-household economy with incomplete insurance and a temporary aggregate income shock. The government observes household characteristics \(h\), but not all permanent income or private savings. A transfer rule \(T(h,s)\ge0\) depends on those observed characteristics and aggregate state \(s\). Let \(\mathcal T\) contain rules satisfying a specified administrative complexity bound, a budget \(E[T(h,s)]\le B(s)\), and a fixed financing scheme. Households optimize consumption and labor, and prices and employment adjust in general equilibrium.
For given nonnegative social weights \(\omega(h)\), utility \(u\), discount factor \(\beta\in(0,1)\), and baseline type distribution \(F\), define
\[
W(T)=E\sum_{t=0}^{H}\beta^t
\int\omega(h)u(c_t(h;T),\ell_t(h;T))\,dF(h).
\]
Determine an implementable maximizing rule within \(\mathcal T\), including the tradeoff between immediate spending responses and insurance against persistent income losses. Identify the consumption-response, income-risk, and reporting parameters needed to rank rules; private information and take-up must enter feasibility rather than being ignored. Compare the optimal conditional rule with a universal transfer under the same budget, financing, and administrative costs. Evaluate performance under plausible misspecification of the type distribution and aggregate shock. The target is a restricted fiscal design problem; a high measured marginal propensity to consume is neither a complete welfare criterion nor sufficient evidence for optimal targeting.

Define observable household characteristics \(h\) on a fixed measurable space and aggregate state \(s\) before the shock. The budget is conditional on each state, \(\int T(h,s)\,dF_s(h)+C_{\rm admin}(T,s)\le B(s)\), with taxes and public debt financing entering equilibrium feasibility. One precise complexity restriction is a fixed partition into \(d\) administrative bins and a bounded nonnegative payment per bin; more general measurable rules need a separate compactness argument. If payments depend only on independently verified characteristics, truthful reporting is not an additional constraint; if characteristics are self-reported, specify report, verification, and penalty rules and impose the resulting participation/incentive constraints. Normalize \(\int\omega(h)\,dF(h)=1\). A maximizing rule may fail to exist, in which case report the supremum and an \(\epsilon\)-optimal implementable rule.
\subsection*{Variants and assumption changes}
The following are \emph{editorial variants}, not additional published open conjectures.
\begin{enumerate}
\item \emph{Verified emergency payments.} Choose bounded transfers to preverified administrative bins for one temporary shock, with fixed tax financing. Compare the welfare optimum with universal payments at equal fiscal and administrative costs.
\item \emph{Automatic stabilization.} Commit before shocks to a state-contingent rule applied repeatedly. Add endogenous labor supply, take-up, and administrative delay, and compare ex ante insurance welfare with a discretionary one-off payment. Self-reporting variants overlap mechanism design.
\end{enumerate}
\subsection*{References and source audit}
\begin{itemize}
\item Adrien Auclert; Matthew Rognlie; Ludwig Straub. \emph{Fiscal and Monetary Policy with Heterogeneous Agents}. 2025. Locator: Section5, Optimal policy, printed pp.21--22; broad policy agenda, not the exact administrative-bin optimization. Primary text checked. \url{https://web.stanford.edu/~aauclert/annual_review_hank.pdf}.
\end{itemize}
Optimal heterogeneous-agent policy is an agenda; exact transfer targeting under administrative constraints is editorial.
\begin{enumerate}\raggedright
\item\hypertarget{definitions:source:30007534:1}{} Adrien Auclert; Matthew Rognlie; Ludwig Straub. 2025. \emph{Fiscal and Monetary Policy with Heterogeneous Agents}. Section5, Optimal policy, printed pp.21--22; broad policy agenda, not the exact administrative-bin optimization.
\url{https://web.stanford.edu/~aauclert/annual_review_hank.pdf}.
DOI: \url{https://doi.org/10.1146/annurev-economics-091624-044646}.
\end{enumerate}

\subsection*{Common conventions: Source mathematical conventions}

These conventions apply to each entry unless explicitly replaced.

\textbf{Models and identification.} A primitive model \(m\) specifies all probability laws, feasible choices, information, equilibrium and measurement rules used by its target. The admissible class \(\mathcal M\) is fixed independently of outcomes. If \(P_O(m)\) is its observable probability law and \(\tau(m)\) the target, its identified set at observable law \(P\) is
\[
 \mathcal I_\tau(P)=\{\tau(m):m\in\mathcal M,\ P_O(m)=P\}.
\]
Point identification means that this set is a singleton. An empty set signals incompatibility of data and assumptions. Coordinate bounds need not describe the joint set. Bounds are sharp only if every retained target is attainable or arbitrarily approachable by compatible models. Finite bounds require explicit restrictions on unobserved quantities; unrestricted confounding, hidden wealth, extrapolation or arbitrary equilibrium selection can yield uninformative sets.

\textbf{Causal interventions.} A potential outcome \(Y_i(p)\) is the outcome under a complete policy regime \(p\), including timing, financing, information and market spillovers. The target population and units are fixed before treatment. With interference, use the complete assignment vector or a justified exposure mapping; individual no-interference notation alone cannot identify an economy-wide intervention. A difference of mean potential outcomes does not require the cross-world joint distribution. Conditioning on post-treatment migration, survival, enrollment or employment changes the estimand and needs separate justification.

\textbf{Statistical statements.} An estimator, confidence set, forecast or test specifies a sampling law, model class and loss/coverage criterion. Uniform \((1-\alpha)\)-coverage means \(\inf_{m\in\mathcal M}\Pr_m\{\tau(m)\in C_n\}\geq1-\alpha\), or an explicitly asymptotic version. The whole parameter space trivially has coverage, so informativeness requires a stated width, risk or power objective. A forecasting improvement is relative to a named benchmark, information set and evaluation population.

\textbf{Equilibrium and optimization.} An equilibrium comprises feasible plans, optimality, beliefs where needed, budget constraints and all market/resource clearing. The primitives or a stated benchmark define each correspondence \(\mathcal E(p;m)\). Nonemptiness is assumed or proved, never inferred just from compact policy sets. A selected-equilibrium target uses a declared selection rule; a robust target quantifies over all permitted equilibria. Use suprema or infima unless attainment is established. An \(\varepsilon\)-optimal policy is feasible and within \(\varepsilon\) of the finite optimum value. Report an empty feasible set separately.

\textbf{Welfare and accounting.} Normative utility, interpersonal normalization, social weights, numeraire, discounting and the use of public receipts are primitives. Behavioral utility need not coincide with the welfare criterion. Household welfare includes ownership income and rents once; adding profits again to compensating variation generally double-counts them. A monetary equivalent is defined by an explicit transfer and price convention, with monotonicity and range conditions. Average effects, mechanism contributions, transfers and welfare losses are different targets. Infinite sums require convergence and dynamic feasible plans require terminal or no-Ponzi conditions.

\textbf{Source versions.} A working-paper date, revision date and journal publication year are distinguished. A DOI for a journal version is not automatically the DOI of an earlier manuscript. Locators refer to the stated version and printed pages where specified. A metadata-only check does not verify a claim about a full-text conclusion. Every entry's source subsection narrows the provenance claim accordingly.
% END ECONOMICS STATEMENT
\end{document}
